\documentclass[10pt,a4paper]{amsart}
\usepackage[margin=19mm]{geometry}
\usepackage{amsmath,amssymb,mathtools}
\usepackage[scr=rsfs,cal=euler]{mathalfa}
\usepackage{xfrac}
\usepackage[unicode,hidelinks,bookmarks=false]{hyperref}
\newcommand{\ie}{i.e.\ }
\newcommand{\cf}{cf.\ }
\newcommand{\resp}{respectively\ }
\newcommand{\Subsection}{\subsection}
\providecommand{\nobreakdash}{\nobreak\hbox{-}}
\setlength{\emergencystretch}{2em}
\multlinegap0pt
\usepackage{mdframed}
\usepackage{mdframed}
\usepackage{mdframed}
\usepackage{mdframed}
\usepackage{amssymb}
\usepackage{tikz}
\usepackage{mdframed}
\usepackage{mathtools}
\usepackage[shortlabels]{enumitem}
\usepackage[normalem]{ulem}
\newtheorem{lemma}{Lemma}
\title{Mesoscopic Rates of Convergence for Complex Wishart Matrices at the Leftmost Spectrum Edge}
\author{Mengchun Cai}
\date{Source: arXiv:2605.12777v1 (CC BY 4.0)}
\begin{document}
\maketitle

\noindent\emph{Source and licence.} Mesoscopic Rates of Convergence for Complex Wishart Matrices at the Leftmost Spectrum Edge. Version arXiv:2605.12777v1 (CC BY 4.0), distributed by arXiv under the Creative Commons Attribution 4.0 International licence, http://creativecommons.org/licenses/by/4.0/, as recorded in the arXiv licence metadata for this exact version and shown on the abstract page https://arxiv.org/abs/2605.12777v1. Full licence notice: \texttt{licenses/arxiv-2605.12777-CC-BY-4.0.txt}.\par\smallskip

\noindent\textit{Editorial scope.} Counted unit: the article's lemma seq (source0.tex lines 1000--1013). The article's complete original proof of it is delivered with it (source0.tex lines 1014--1069, one \texttt{proof} environment of 56 lines). No mathematics is written, repaired or re-derived: the statement and the proof are the article's own lines, in the article's own notation, and no retained line is substituted or re-flowed.  No further source-local context is needed: every cross-reference of every retained line resolves inside the two delivered spans.  Nothing was omitted from any retained span, so the package carries no omission record.  No retained line contains a citation, so the wrapper carries no bibliography and no bibliography entry.  No retained line contains a figure environment, an image-inclusion command or drawing code, so no image is delivered and the package declares no asset dependency.  Editorial formatting only, in addition: an AMS \texttt{amsart} preamble that restates the article's own declarations and macros used by the retained lines, namely the environments \texttt{lemma} on separate counters, together with the standard packages those lines need.  The retained environments print in this document as Lemma 1 in order of appearance, whereas the article prints the same items with the numbering of its own full document.  The complete native source of the article as distributed in the arXiv e-print for this exact version is bundled unchanged as \texttt{sources/200451/source0.tex}.
\begin{lemma}\label{seq}
    If $z=z_1-\epsilon$ for some suitable $\epsilon$ (say less or equal to $\frac{\mu_{n,N}}{2\kappa_{N}}$), then there exists a sequence $s_N$ with finite limit such that
    \begin{equation}\label{a4zeta}
    \zeta^{\frac{3}{2}}(z)=\frac{\sqrt{\epsilon^3(z_{2}-z_{1})}}{2z_{1}}\left(1+\frac{3}{5}s_{N}\epsilon+O(\epsilon^2)\right)
\end{equation}
and hence,
\begin{equation}\label{a4zeta2}
   \kappa_{N}^{\frac{2}{3}}\zeta(z)=\frac{\epsilon\kappa_{N}}{\sigma_{n,N}}\left(1+\frac{2\epsilon s_{N}}{5}+O(\epsilon^2)\right),
\end{equation}
\begin{equation}\label{a4inverse}
    \left(\frac{\kappa_{N}}{\sigma^3_{n,N}}\right)^{\frac{1}{6}}\tilde{f}^{-\frac{1}{4}}(z)=1-\frac{2}{5} s_{N}\epsilon+O(\epsilon^2).
\end{equation}
Here $z_1$, $z_2$, $\kappa$ and $s$ are all double-indexed sequences with respect to $(j,k)\in\mathbb{N}^2$. We will abbreviate them as in this lemma if $(j,k)=(n,N)$. Moreover, the L--G transform $\zeta$ and induced $\tilde{f}$ is always assumed to be around $z_1$ (with parameter $(n,N)$ here).
\end{lemma}

\begin{proof}
   According to the definition, 
\begin{equation*}
        \frac{2}{3}\zeta^{\frac{3}{2}}(z)=\frac{2}{3}\zeta^{\frac{3}{2}}\left(\frac{x_N(s)}{\kappa_N}\right)=\int_{z_1-\epsilon}^{z_1}\frac{\sqrt{(t-z_1)(t-z_2)}}{2t}dt.
\end{equation*}
Recall $\frac{\kappa_N}{\sigma_{n,N}^3}=\frac{z_2-z_1}{4z_1^2}$ by our construction. Define $A_N:=z_2-z_1=2\sqrt{4-\omega_N^2}$ and substitute $t$ by $y:=\frac{z_1-t}{\epsilon}$, we have:
\begin{equation*}
          \frac{2}{3}\zeta^{\frac{3}{2}}(z)=\frac{\sqrt{\epsilon^3A_N}}{2z_1}\int_0^1\sqrt y\frac{\sqrt{1+\frac{\epsilon}{A_N}y}}{1-\frac{\epsilon}{z_1}y}dy.
\end{equation*}
Since both $A_N$ and $z_1$ have finite nonzero limits as $\gamma\in(0,1)$, $\frac{\epsilon y}{A_N}$ and $\frac{\epsilon y}{z_1}$ are both small for any $y\in[0,1]$ provided $\epsilon$ in our region. The dominated convergence theorem allows us to perform the Taylor expansion inside the integral as:
\begin{equation}\label{int4epsilon}
\begin{split}
    \int_0^1\sqrt y\frac{\sqrt{1+\frac{\epsilon}{A_N}y}}{1-\frac{\epsilon}{z_1}y}dt&=\int_0^1\sqrt{y}\left(1+\epsilon y(\frac{1}{z_1}+\frac{1}{2A_N})+O(\epsilon^2)\right)dy\\
    &=\frac{2}{3}+\frac{2\epsilon}{5}\left(\frac{1}{z_1}+\frac{1}{2A_N}\right)+O(\epsilon^2).
\end{split}
\end{equation}
In fact, since
\begin{equation*}
    \lim_{N\rightarrow\infty}z_1=2-\frac{4\sqrt\gamma}{1+\gamma}=\frac{2(1-\sqrt\gamma)^2}{1+\gamma}
\end{equation*}
and
\begin{equation*}
    \lim_{N\rightarrow\infty}\frac{\mu_{n,N}}{\kappa_N}=\lim_{N\rightarrow\infty}\frac{2(\sqrt n-\sqrt N)^2}{N+n+1}=\frac{2(1-\sqrt\gamma)^2}{1+\gamma},
\end{equation*}
 we know $\epsilon\le\frac{\mu_{n,N}}{2\kappa_N}<\frac{2z_1}{3}$ for $N$ large enough by our restriction for $\epsilon$ and hence, $1-\frac{\epsilon}{z_1}y\ge \frac{1}{3}$ for any $y\in[0,1]$. Therefore, for $N$ large, the integrand in \eqref{int4epsilon} is bounded by $3\sqrt y\sqrt{1+C(\gamma)y}$, which is in $L^1[0,1]$. That means the dominated convergence theorem applies.

Define $s_N:=\frac{1}{z_1}+\frac{1}{2A_N}$. We know $s_N$ has finite limit and
\begin{equation*}
    \zeta^{\frac{3}{2}}(z)=\frac{\sqrt{\epsilon^3A_N}}{2z_1}\left(1+\frac{3}{5}s_N\epsilon+O(\epsilon^2)\right),
\end{equation*}
which is exactly \eqref{a4zeta}.

Recall $\frac{A_N}{4z_1^2}=\zeta'(z_1)=\frac{\kappa_N}{\sigma_{n,N}^3}$ by our construction, from \eqref{a4zeta},
\begin{equation*}
    \begin{split}
        \kappa_N^{\frac{2}{3}}\zeta&=\kappa_N^{\frac{2}{3}}\epsilon\left(\frac{A_N}{4z_1^2}\right)^\frac{1}{3}\left(1+\frac{2s_N\epsilon}{5}+O(\epsilon^2)\right)\\
        &=\frac{\epsilon\kappa_N}{\sigma_{n,N}}\left(1+\frac{2s_N\epsilon}{5}+O(\epsilon^2)\right)
    \end{split}.
\end{equation*}
Now write $f$ as:
\begin{equation*}
    f(z_1-\epsilon)=\frac{\epsilon(\epsilon+(z_2-z_1))}{4z_1^2(1-\frac{\epsilon}{z_1})^2}=\frac{A_N\epsilon(1+\frac{\epsilon}{A_N})}{4z_1^2(1-\frac{\epsilon}{z_1})^2}.
\end{equation*}
Again, performing the Taylor expansion for small $\epsilon$ leads to
\begin{equation*}
    f^{-\frac{3}{2}}(z)=\left(\frac{A_N\epsilon}{4z_1^2}\right)^{-\frac{3}{2}}\left(1-3s_N\epsilon+O(\epsilon^2)\right).
\end{equation*}
Recall $\tilde{f}=\frac{f}{\zeta}$, combing the previous results yields:
\begin{equation*}
    \tilde{f}^{-\frac{1}{4}}(z)=\left(\zeta^{\frac{3}{2}}(z)f^{-\frac{3}{2}}(z)\right)^{\frac{1}{6}}=\left(\frac{4z_1^2}{A_N}\right)^{\frac{1}{6}}\left(1-\frac{2}{5}s_N\epsilon+O(\epsilon^2)\right)
\end{equation*}
and hence,
\begin{equation*}
    \left(\frac{\kappa_N}{\sigma^3_{n,N}}\right)^{\frac{1}{6}}\tilde{f}^{-\frac{1}{4}}(z)=\left(\frac{\kappa_N}{\sigma^3_{n,N}}\times\frac{4z_1^2}{A_N}\right)^{\frac{1}{6}}\left(1-\frac{2}{5}s_N\epsilon+O(\epsilon^2)\right)=1-\frac{2}{5}s_N\epsilon+O(\epsilon^2)
\end{equation*}
for $N$ large enough.\end{proof}

\end{document}
